% ---------------------------------------------------------------
%  Section VIII -- Discussion
% ---------------------------------------------------------------

The core finding of this paper is architectural: the IRIDIUM platform
demonstrates that a real-data urban mobility stack can be constructed
with explicit, machine-readable provenance and data-status semantics,
without requiring access to a complete set of operational feeds at the
time of development or initial deployment.
Three design decisions are central to this property, and each has
broader implications for the design of transport data platforms in
data-scarce environments.

\textbf{Strict no-synthetic-substitution policy.}
When a source is absent or unreachable, the assembler returns empty
fields with a specific status code rather than populating the response
with simulated data.
This prevents the failure mode in which a system appears to perform well
during testing because a synthetic fallback is indistinguishable from
live data.
It also means that the system is directly usable in production with
partial data: a planner who queries the API receives an honest representation
of what is known and what is not, enabling informed interpretation rather
than silent overconfidence.

\textbf{Composable, machine-readable provenance records.}
The provenance list $\Pi$ in each API response (Eq.~\ref{eq:provenance})
is versioned and structured, enabling automatic audit trails and
incremental data-quality monitoring.
As new sources are activated under data-sharing agreements, their
contributions appear in $\Pi$ without requiring changes to the API schema.
This composability property reduces the cost of incremental platform
maturation, which is the realistic mode of development for urban data
platforms in the Global South.

\textbf{Architectural separation of concerns.}
The analytics layer depends on the twin snapshot interface, not on
specific data backends.
Substituting PostgreSQL/PostGIS with an alternative graph store,
replacing the heuristic forecaster with a trained ST-GNN, or integrating
OpenTripPlanner for multimodal routing, requires no modifications to
the routing, equity, or anomaly modules.
This independence reduces the cost of each individual upgrade and limits
the blast radius of integration failures.

\textbf{Interpretation of weather results.}
The persistent Baku--Quba temperature differential of
$\overline{\Delta T} = 2.50\,\text{\textdegree{}C}$ reflects the expected
superposition of orographic cooling ($\approx 3.8\,\text{\textdegree{}C}$
at the standard lapse rate) and Caspian thermal buffering of Baku.
This validates the multi-city API call structure and confirms that the
Open-Meteo integration is producing physically coherent output.
Operationally, it implies that weather-conditioned features and
adverse-weather thresholds for anomaly severity scoring must be
city-specific throughout Azerbaijan, not shared across regions with
differing elevation profiles.
The precipitation analysis confirms that live weather data can drive
$w_t$ in Eq.~(\ref{eq:anomaly_severity}) in real-time, providing
a weather-aware dimension to the anomaly severity score that is absent
from purely traffic-speed-based detectors.

\textbf{Comparison with related platforms.}
Table~\ref{tab:related_systems} (Section~\ref{sec:related}) shows that
existing ST-GNN benchmarks (DCRNN, Graph WaveNet, ASTGCN) operate on
curated traffic datasets from mature data ecosystems that are not yet
matched by the baseline Azerbaijani provider configuration described here.
IRIDIUM provides the architectural scaffold on which analogous evaluations
can be conducted once operator agreements are in place, without requiring
a redesign of the analytics layer.
The provenance protocol and the explicit equity-analysis layer are not
standard components of the benchmark systems reviewed, while the federated
learning pathway remains uncommon in urban mobility platform papers.

\textbf{Significance for Azerbaijani cities.}
At the current development stage, the practical value of IRIDIUM for
Azerbaijani urban planning lies in three areas.
First, the documented data stack provides a reference architecture for
city planners who need to understand exactly which data agreements are
required and how each feed would be used.
Second, the OSM/PostGIS pipeline is immediately operational for network
analysis, accessibility queries, and single-modal routing once credentials
are obtained.
Third, the Open-Meteo integration provides a live atmospheric input for
anomaly severity scoring and route quality degradation warnings without
any additional operator agreements, which is immediately deployable.

\textbf{Outstanding uncertainties.}
Forecasting accuracy under Azerbaijani traffic dynamics, routing quality
under multimodal load, equity score calibration, and anomaly detection
performance against a labeled incident log remain open.
These are operational constraints of the current deployment stage,
not architectural limitations; each has a resolution path detailed in
Section~\ref{sec:conclusion}.
